\documentclass[11pt,a4paper]{article}
\usepackage[margin=25mm]{geometry}
\usepackage{fontspec}
\setmainfont{TeX Gyre Pagella}
\usepackage{amsmath,amssymb,mathtools}
\usepackage{graphicx}
\usepackage{xcolor}
\usepackage[unicode,colorlinks=true,linkcolor=blue,urlcolor=blue]{hyperref}
\usepackage{bookmark}
\usepackage{enumitem}
\setlist{itemsep=3pt,topsep=5pt}
\setlength{\emergencystretch}{3em}
\newcommand{\tightlist}{\setlength{\itemsep}{0pt}\setlength{\parskip}{0pt}}
\newcommand{\passthrough}[1]{#1}
\newcommand{\pandocbounded}[1]{#1}
\makeatletter
\def\maxwidth{\ifdim\Gin@nat@width>\linewidth\linewidth\else\Gin@nat@width\fi}
\def\maxheight{\ifdim\Gin@nat@height>0.75\textheight0.75\textheight\else\Gin@nat@height\fi}
\makeatother
\setkeys{Gin}{width=\maxwidth,height=\maxheight,keepaspectratio}
\hypersetup{pdfauthor={OpenAI GPT-6.1 Sol, Codex, Ultra effort}}
\begin{document}
\section{Lubin--Tate division fields}\label{lubintate-division-fields}

\emph{Written by OpenAI GPT-6.1 Sol in Codex, Ultra effort, October
2026. Self-checked by the writing AI; no independent review is claimed.
Public domain (CC0).}

A Lubin--Tate formal module gives a way to divide by a uniformizer. The
points killed by \(\pi^n\) form one cyclic module over
\(\mathcal O_K/\pi^n\). Its generators generate a field, and the units
of that quotient act as every automorphism of the field. We prove these
assertions by counting polynomial roots and applying Eisenstein's
criterion.

Fix a nonarchimedean local field \(K\), its valuation ring
\(\mathcal O\), a uniformizer \(\pi\), and residue field
\(\mathbf F_q\). Normalize \(v_K(\pi)=1\) and extend this valuation to a
separable closure \(K^s\). For a Lubin--Tate series
\(f\in\mathcal F_\pi\), write \(F_f\) and \([a]_f\) for the formal group
and scalar action constructed in
\href{../formal-groups-and-lubin-tate-modules.html}{Formal groups and
Lubin--Tate modules}, Theorem 7.3.

We also use
\href{https://kokunoyumeto.github.io/open-math-courses/courses/NT-CFT/proof-dependencies.html\#NT-LOC-07}{Unramified
and totally ramified extensions}, Theorem 5.1: a root of a monic
Eisenstein polynomial of degree \(d\) over a complete discretely valued
field generates a totally ramified extension of degree \(d\), and that
root is a uniformizer. Finite extensions are complete with the uniquely
extended valuation, by
\href{https://kokunoyumeto.github.io/open-math-courses/courses/NT-CFT/proof-dependencies.html\#NT-LOC-04}{Extensions
of complete valued fields}, Theorem 1.2.

\textbf{Prerequisite proof availability.} The named results below
identify specific programme lessons. The
\href{https://kokunoyumeto.github.io/open-math-courses/courses/NT-CFT/proof-dependencies.html\#lesson-8}{prerequisite
record} distinguishes published proofs, supplied owner texts awaiting
publication, and missing full proofs. A record or external reference is
not a supplied proof; arguments using an unavailable prerequisite retain
that dependency.

\subsection{1. Evaluating the formal
module}\label{evaluating-the-formal-module}

Put \[
\mathfrak m_s=\{x\in K^s:v_K(x)>0\}.
\] Here \(0\) has valuation \(+\infty\). For finitely many elements of
\(\mathfrak m_s\), choose a finite extension of \(K\) containing them.
Integral power series converge there, since their terms of degree
tending to infinity have valuation tending to infinity. Thus \(F_f\),
its inverse and every \([a]_f\) can be evaluated on \(\mathfrak m_s\).
This does not require the separable closure itself to be complete.

The identities proved in lesson 7 give \(\mathfrak m_s\) an
\(\mathcal O\)-module structure, with addition \(F_f\). Define \[
T_n(f)=F_f[\pi^n]
 =\{x\in\mathfrak m_s:f^{\circ n}(x)=0\},\qquad T_0(f)=\{0\}.
\tag{1}
\] The superscript \(\circ n\) denotes composition, not the \(n\)th
power of a series. A point is \textbf{primitive at level \(n\)} if it
belongs to \(T_n(f)\setminus T_{n-1}(f)\).

The particularly useful series \[
f_0(X)=\pi X+X^q
\tag{2}
\] is a polynomial. It permits a finite root calculation. For an
arbitrary \(f\), Theorem 7.3 supplies an integral isomorphism \[
h=[1]_{f,f_0}:F_{f_0}\longrightarrow F_f,\qquad
h(X)=X+\text{higher terms},\qquad f\circ h=h\circ f_0.
\tag{3}
\] Its inverse is integral, and both series can be evaluated in every
finite extension on its maximal ideal. Consequently \(h\) identifies
\(T_n(f_0)\) with \(T_n(f)\), including their scalar actions and
primitive points.

\subsubsection{Proposition 8.1. The division
module}\label{proposition-8.1.-the-division-module}

For each \(n\geq1\), \(T_n(f)\) has \(q^n\) elements and is a free
module of rank one over \(\mathcal O/\pi^n\mathcal O\). If \(\lambda\)
is any primitive point, the isomorphism is \[
\mathcal O/\pi^n\mathcal O\longrightarrow T_n(f),\qquad
a\longmapsto[a]_f(\lambda).
\tag{4}
\]

\textbf{Proof.} First use \(f_0\). The polynomial \(f_0^{\circ n}\) is
monic of degree \(q^n\). All its roots have positive valuation. Indeed,
for \(v_K(x)<0\), \[
v_K(f_0(x))=qv_K(x)<0,
\] because the two summands have distinct valuations \(1+v_K(x)\) and
\(qv_K(x)\). For \(v_K(x)=0\), its value likewise has valuation zero.
Neither type of point can reach zero under iteration.

At every point of positive valuation, \[
f_0'(x)=\pi+q x^{q-1}\ne0.
\tag{5}
\] In characteristic \(p\), the second summand is zero. In
characteristic zero, \(v_K(q)\geq1\), so that summand has valuation
strictly greater than 1, whereas \(\pi\) has valuation 1. This also
gives \(f_0'(0)=\pi\ne0\). The derivative of the iterate is the product
of these nonzero derivatives along the orbit of a root. Thus all \(q^n\)
roots are simple and lie in \(K^s\). We have \(|T_n(f_0)|=q^n\).

The nesting \(T_{n-1}\subset T_n\) is strict, so primitive points exist.
Fix one, \(\lambda\). The kernel of \(a\mapsto[a]_{f_0}(\lambda)\)
contains \(\pi^n\mathcal O\). Every nonzero \(a\in\mathcal O\) has the
form \(\pi^r u\) with \(u\) a unit. The series \([u]_{f_0}\) is an
invertible module map. If \(r<n\), then \([\pi^r]_{f_0}(\lambda)\ne0\),
since otherwise \(\lambda\) would be killed by \(\pi^{n-1}\). Hence
\([a]_{f_0}(\lambda)\ne0\). The kernel is exactly \(\pi^n\mathcal O\).

The resulting injection from \(\mathcal O/\pi^n\mathcal O\), a set of
\(q^n\) elements, onto a subset of \(T_n(f_0)\) is a bijection. Finally
(3) transports the count and this module isomorphism to every \(f\).
\(\square\)

In particular, the primitive points correspond exactly to the unit
classes. Their number is \[
D_n=|(\mathcal O/\pi^n\mathcal O)^\times|
      =(q-1)q^{n-1}.
\tag{6}
\]

\subsection{2. The field and its
automorphisms}\label{the-field-and-its-automorphisms}

Define \[
K_{\pi,n}=K(T_n(f)).
\tag{7}
\] This notation will be justified by proving independence from \(f\);
the uniformizer \(\pi\) remains part of the notation.

\subsubsection{Proposition 8.2. Independence of the Lubin--Tate
series}\label{proposition-8.2.-independence-of-the-lubintate-series}

The field in (7) is the same for every \(f\in\mathcal F_\pi\). For any
primitive \(\lambda\in T_n(f)\), \[
K_{\pi,n}=K(\lambda).
\tag{8}
\]

\textbf{Proof.} An integral series with coefficients in \(K\), evaluated
at a point of positive valuation in a finite extension \(L/K\), has its
value in \(L\), by completeness. Thus (3) maps every point
\(x\in T_n(f_0)\) into \(K(x)\). Its integral inverse maps \(h(x)\) back
into \(K(h(x))\). Therefore \(K(x)=K(h(x))\), and the fields generated
by all points are equal.

By Proposition 8.1, every element of \(T_n(f)\) is \([a]_f(\lambda)\)
for some \(a\in\mathcal O\). These values lie in the complete finite
field \(K(\lambda)\). It contains all the points, proving (8).
\(\square\)

If \(\sigma\in\operatorname{Gal}(K^s/K)\), then \[
\sigma([a]_f(x))=[a]_f(\sigma x),\qquad
\sigma(F_f(x,y))=F_f(\sigma x,\sigma y).
\tag{9}
\] To justify passing through an infinite series, put the inputs and
their conjugates in a finite Galois extension. Its automorphisms
preserve the uniquely extended valuation and are continuous; each
coefficient is fixed by \(\sigma\), so the identities hold for finite
truncations and then for their limits.

For \(f_0\), \(K_{\pi,n}\) is the splitting field of the separable
polynomial \(f_0^{\circ n}\), hence is finite Galois. Proposition 8.2
gives the same conclusion for every \(f\). Equation (9) defines an
injection \[
\operatorname{Gal}(K_{\pi,n}/K)
 \hookrightarrow\operatorname{Aut}_{\mathcal O}(T_n(f))
 \simeq(\mathcal O/\pi^n\mathcal O)^\times.
\tag{10}
\] It is injective because the division points generate the field. An
automorphism of the cyclic module is multiplication by a unique unit
class.

\subsubsection{Theorem 8.3. Degree, ramification and Galois
group}\label{theorem-8.3.-degree-ramification-and-galois-group}

The extension \(K_{\pi,n}/K\) is totally ramified of degree \(D_n\),
every primitive point is a uniformizer of this extension, and (10) is an
isomorphism.

\textbf{Proof.} Continue first with \(f_0\). Put
\(A_n(X)=f_0^{\circ(n-1)}(X)\), where \(A_1(X)=X\). Polynomial
composition gives the exact factorization \[
f_0^{\circ n}(X)=A_n(X)\bigl(\pi+A_n(X)^{q-1}\bigr).
\] Thus \[
P_n(X)=\frac{f_0^{\circ n}(X)}{f_0^{\circ(n-1)}(X)}
       =\pi+A_n(X)^{q-1}
\tag{11}
\] is a monic polynomial of degree \(D_n\). Since \(A_n(0)=0\), its
constant coefficient is \(\pi\). Reduction modulo \(\pi\) gives
\(A_n(X)\equiv X^{q^{n-1}}\), and therefore \(P_n(X)\equiv X^{D_n}\).
Every other nonleading coefficient is divisible by \(\pi\), while its
constant coefficient is not divisible by \(\pi^2\). This is precisely
the Eisenstein condition.

A primitive \(\lambda\) is a root of \(P_n\), because
\(A_n(\lambda)\ne0\) while \(f_0^{\circ n}(\lambda)=0\). Here the degree
and ramification consequences of the Eisenstein condition can also be
verified directly. Write \(d=D_n\) and \(t=v_K(\lambda)>0\). In
\(P_n(\lambda)\), the leading term has valuation \(dt\), the constant
term has valuation 1, and every intervening term has valuation at least
\(1+it>1\). A sum equal to zero cannot have a unique term of least
valuation: after dividing by that term, reduction in the residue field
would give \(1=0\). Therefore \(dt=1\). If \(e\) is the ramification
index of \(K(\lambda)/K\), its value group in this normalization is
\(e^{-1}\mathbf Z\); thus \(d\) divides \(e\). The degree formula for
complete discretely valued fields gives \(d\leq e\leq[K(\lambda):K]\),
whereas the polynomial \(P_n\) gives the reverse bound
\([K(\lambda):K]\leq d\). All these quantities are consequently \(d\),
the residue degree is 1, and \(\lambda\) has normalized valuation 1 in
its field. In particular \[
[K(\lambda):K]=D_n,
\] total ramification, and the uniformizer assertion. By (8), this is
the entire division field. Its Galois group and the right side of (10)
both have \(D_n\) elements; the injection is consequently onto.

For general \(f\), the strict isomorphism (3) preserves the generated
field and the valuation of each nonzero point: if \(v_K(x)>0\), every
higher term of \(h(x)-x\) has valuation strictly greater than
\(v_K(x)\). A primitive \(h(\lambda)\) therefore remains a uniformizer.
This proves all the assertions. \(\square\)

The polynomial argument uses the specific choice (2). An arbitrary
Lubin--Tate series can have infinitely many nonzero coefficients; it
need not supply a polynomial quotient as in (11). The integral
isomorphism transfers the field and module conclusions without making
such an assumption.

\subsubsection{Proposition 8.4. A uniformizer is a
norm}\label{proposition-8.4.-a-uniformizer-is-a-norm}

For every \(n\geq1\), \[
\pi\in N_{K_{\pi,n}/K}(K_{\pi,n}^{\times}).
\tag{12}
\]

\textbf{Proof.} Choose a primitive \(\lambda\) for \(f_0\). Its minimal
polynomial is \(P_n\), with constant coefficient \(\pi\), so \[
N_{K_{\pi,n}/K}(\lambda)=(-1)^{D_n}\pi.
\] Since \(N(-1)=(-1)^{D_n}\), we obtain \[
N_{K_{\pi,n}/K}(-\lambda)=\pi.
\tag{13}
\] The field is independent of the chosen series, so (12) holds for it
in every presentation. \(\square\)

When \(q\) is even and \(n=1\), \(D_1=q-1\) is odd, and
\(N(\lambda)=-\pi\). Formula (13) includes this case. A primitive point
for a different series is still a uniformizer, but its norm need not
have the particular value of the polynomial coordinate just used.

\subsection{3. The infinite tower}\label{the-infinite-tower}

The nesting of the division modules gives fields \[
K_{\pi,1}\subset K_{\pi,2}\subset\cdots,\qquad
K_\pi=\bigcup_{n\geq1}K_{\pi,n}.
\tag{14}
\]

\subsubsection{Corollary 8.5. The Galois group of the
tower}\label{corollary-8.5.-the-galois-group-of-the-tower}

There is a canonical isomorphism of topological groups, characterized by
its action on division points, \[
\operatorname{Gal}(K_\pi/K)\simeq\mathcal O^\times.
\tag{15}
\]

\textbf{Proof.} Multiplication by \(\pi\) maps \(T_{n+1}(f)\) onto
\(T_n(f)\): in the cyclic module \(\mathcal O/\pi^{n+1}\), its image has
\(q^n\) elements, lies in \(T_n\), and hence equals it. Choose primitive
points compatibly with \([\pi]_f(\lambda_{n+1})=\lambda_n\). A preimage
of a primitive level-\(n\) point must be primitive at level \(n+1\).

If \(\sigma(\lambda_{n+1})=[u]_{f}(\lambda_{n+1})\), applying
\([\pi]_f\) shows \(\sigma(\lambda_n)=[u]_f(\lambda_n)\). Thus the
restriction maps in Theorem 8.3 correspond to reduction of unit classes.
Infinite Galois theory from lesson 1 now yields \[
\operatorname{Gal}(K_\pi/K)
 \simeq\varprojlim_n(\mathcal O/\pi^n\mathcal O)^\times.
\]

For completeness, this last inverse limit is \(\mathcal O^\times\).
Given compatible classes, choose representatives
\(u_n\in\mathcal O^\times\). They satisfy
\(u_{n+1}-u_n\in\pi^n\mathcal O\), so form a Cauchy sequence and have a
limit \(u\in\mathcal O\). Its nonzero residue makes it a unit, and it
represents all the classes. Uniqueness follows from
\(\bigcap_n\pi^n\mathcal O=\{0\}\). Reduction is continuous, and its
kernels \(1+\pi^n\mathcal O\) form a neighborhood basis, so the
bijection and its inverse are continuous. \(\square\)

The isomorphism uses scalar action and is independent of the chosen
compatible generators. All finite levels are abelian, as is the tower.

\subsection{4. Two concrete
descriptions}\label{two-concrete-descriptions}

For \(K=\mathbf Q_p\), \(\pi=p\), choose \[
f(X)=(1+X)^p-1,\qquad F_f(X,Y)=X+Y+XY.
\] Lesson 7 proves \([a]_f(X)=(1+X)^a-1\) for \(a\in\mathbf Z_p\).
Iterating gives \(f^{\circ n}(X)=(1+X)^{p^n}-1\). Therefore \[
T_n(f)=\{\zeta-1:\zeta^{p^n}=1\},\qquad
K_{p,n}=\mathbf Q_p(\zeta_{p^n}).
\tag{16}
\] The root count and positive-valuation assertion can also be obtained
by transferring from \(f_0\), as above. For a primitive \(p^n\)th root
\(\zeta\), (10) says precisely \[
\sigma(\zeta)=\zeta^u,\qquad u\in(\mathbf Z/p^n\mathbf Z)^\times.
\tag{17}
\] For a \(p\)-adic scalar, choose an ordinary integer congruent to it
modulo \(p^n\); Proposition 8.1 makes their actions identical on
\(T_n\). Thus the exponent in (17) has an unambiguous meaning.

For the polynomial \(f_0\) over general \(K\), its nonzero first-level
points satisfy \[
\lambda^{q-1}=-\pi,\qquad
K_{\pi,1}=K\bigl((-\pi)^{1/(q-1)}\bigr).
\tag{18}
\] This is a totally ramified extension of degree \(q-1\). If \(q=2\),
its degree is one and \(\lambda=-\pi\) already belongs to \(K\).

Changing the series for a fixed uniformizer preserves the tower.
Changing the uniformizer can change even its first field. For example,
over \(\mathbf Q_3\), the choices \(3\) and \(-3\) yield
\(\mathbf Q_3(\sqrt{-3})\) and \(\mathbf Q_3(\sqrt3)\); Exercise 4
proves they are different. The next lesson will explain why adjoining
the entire unramified tower removes this dependence.

\subsection{5. Exercises and complete
solutions}\label{exercises-and-complete-solutions}

\subsubsection{Exercise 1 --- easy}\label{exercise-1-easy}

For \(K=\mathbf Q_p\), verify the torsion and Galois descriptions
(16)--(17) using the multiplicative formal group.

\textbf{Solution.} Composition of \(f(X)=(1+X)^p-1\) multiplies the
exponent by \(p\), giving \(f^{\circ n}(X)=(1+X)^{p^n}-1\) by induction.
Its zeros are exactly \(\zeta-1\). They lie in the positive-valuation
domain because this torsion set is the image under (3) of the polynomial
torsion set. A primitive point corresponds to a primitive root:
otherwise its exponent order would divide \(p^{n-1}\). All other roots
are powers of it, so the generated field is
\(\mathbf Q_p(\zeta_{p^n})\).

For integer \(a\), \([a]_f(\zeta-1)=\zeta^a-1\) follows from the
multiplicative law. For any \(a\in\mathbf Z_p\), choose that integer
modulo \(p^n\); their difference kills the division module. Theorem 8.3
makes every unit class a Galois automorphism, giving (17) and degree
\((p-1)p^{n-1}\), including \(p=2\).

\subsubsection{Exercise 2 --- medium}\label{exercise-2-medium}

Prove that the Galois action on \(T_n(f)\) commutes with every scalar
\([a]_f\).

\textbf{Solution.} Put a division point and its conjugate in a finite
Galois extension \(L/K\). The coefficients of \([a]_f\) belong to
\(\mathcal O\subset K\). For each finite truncation, applying \(\sigma\)
to its value equals evaluating it at \(\sigma x\). The unique extended
valuation on \(L\) is \(\sigma\)-invariant, so \(\sigma\) is an isometry
and preserves limits. Passing to the convergent series gives (9).
Applying the same argument to \(F_f\) proves additivity. In particular
\(\sigma\) takes points killed by \(\pi^n\) to points killed by
\(\pi^n\), and is an \(\mathcal O\)-module automorphism there.

\subsubsection{Exercise 3 --- medium}\label{exercise-3-medium}

For the polynomial choice \(f_0=\pi X+X^q\), prove that
\(f_0^{\circ n}/f_0^{\circ(n-1)}\) is Eisenstein. Explain the role of
this choice for a general series.

\textbf{Solution.} With \(A=f_0^{\circ(n-1)}\), composition is
\(f_0(A)=\pi A+A^q\). The quotient is the polynomial \(\pi+A^{q-1}\).
The polynomial \(A\) is monic of degree \(q^{n-1}\), has constant
coefficient zero, and reduces to \(X^{q^{n-1}}\) modulo \(\pi\). Hence
the quotient is monic of degree \((q-1)q^{n-1}\), has constant
coefficient exactly \(\pi\), and reduces to the pure power
\(X^{(q-1)q^{n-1}}\). These are all the Eisenstein conditions.

For a general power series in \(\mathcal F_\pi\), its iterate is not
necessarily a polynomial. The assertions about torsion and fields follow
instead from the integral strict isomorphism with \(f_0\), which
preserves primitive points, their fields and their valuations. This
supplies the same degree and ramification conclusion without assuming
finite polynomial degree for an infinite series.

\subsubsection{Exercise 4 --- hard}\label{exercise-4-hard}

Let \(K=\mathbf Q_p\) with \(p\) odd. For uniformizers \(\pi,\pi'\),
prove \[
K_{\pi,1}=K_{\pi',1}
\quad\Longleftrightarrow\quad
\pi'/\pi\equiv1\pmod p.
\tag{19}
\]

\textbf{Solution.} Put \(m=p-1\), \(u=\pi'/\pi\), and choose
\(\alpha^m=-\pi\), \(\beta^m=-\pi'\). Equation (18) gives the two
fields. If they are equal to \(L\), both \(\alpha,\beta\) are
uniformizers of the same totally ramified extension. Their ratio is a
unit of \(L\), and \[
(\beta/\alpha)^m=u.
\] The residue field of \(L\) is \(\mathbf F_p\). Every nonzero residue
has \(m\)th power 1, so reducing this equality gives
\(u\equiv1\pmod p\).

Conversely suppose \(u\in1+p\mathbf Z_p\). We construct
\(a\in1+p\mathbf Z_p\) with \(a^m=u\). Start \(a_0=1\), and put \[
a_{j+1}=a_j+\frac{u-a_j^m}{m a_j^{m-1}}.
\] The denominator is a \(p\)-adic unit. If \(a_j\in1+p\mathbf Z_p\) and
the error \(u-a_j^m\) has valuation at least \(r\geq1\), the increment
has valuation at least \(r\). Expanding the \(m\)th power cancels its
linear error; all remaining terms contain at least the square of the
increment, and have integral coefficients. The next error has valuation
at least \(2r\). Thus the increments tend to zero with valuations at
least \(2^j\), the sequence converges in \(1+p\mathbf Z_p\), and its
limit \(a\) satisfies \(a^m=u\). Now \(a\alpha\) is a root of
\(X^m+\pi'\). It generates \(K(\alpha)\), since \(a\in K^\times\); hence
the two fields are equal. This proves (19).

\subsection{Editable edition}\label{editable-edition}

The reading edition provides the complete LaTeX source of this lesson,
the cumulative course LaTeX and the editable source ZIP. The archive
contains all twenty-four Markdown lessons, complete LaTeX bodies,
original diagrams, metadata and reproduction instructions.

\subsection{References}\label{references}

Sections 1--2 prove root counts, the division module, independence of
the defining series and the full Galois group. The valuation calculation
in Theorem 8.3 proves the degree and total ramification directly. The
norm sign and the whole tower are retained.

\begin{itemize}
\tightlist
\item
  \href{https://www.jmilne.org/math/CourseNotes/CFT.pdf}{J. S. Milne,
  Class Field Theory, version 4.03}.
\item
  \href{https://arxiv.org/abs/math/0606108v2}{Teruyoshi Yoshida, Local
  class field theory via Lubin--Tate theory, arXiv:math/0606108v2}.
\end{itemize}

The \href{../free-proofs.html}{proof guide} gives the lesson sequence
and the exact prerequisite record. External references accompany the
written arguments.

\end{document}
